% ==================================================
% Summation Notation
% Discrete Calculus — Foundations
% ==================================================

\section{Summation Notation}

\begin{definitionbox}{Definition (Summation)}
Let \(f\) be a function defined on the integers.
For integers \(a \le b\), the expression
\[
\sum_{k=a}^{b} f(k)
\]
denotes the sum \(f(a) + f(a+1) + f(a+2) + \cdots + f(b)\).
The symbol \(\sum\) (Greek capital sigma) represents summation;
\(k\) is the \emph{index of summation}; \(a\) and \(b\) are the
\emph{lower} and \emph{upper limits}; \(f(k)\) is the \emph{summand}.
\end{definitionbox}
\label{def:summation}

\begin{remark*}[Fully quantified form]
\(\displaystyle\sum_{k=a}^{b} f(k) \;:=\; f(a) + f(a+1) + \cdots + f(b)\),
where \(a,b \in \mathbb{Z}\), \(a \le b\), and \(f : \mathbb{Z}\to\mathbb{R}\).
\end{remark*}

\begin{remark*}[Dummy variable]
The index of summation is a \emph{dummy variable}:
\(\displaystyle\sum_{k=a}^{b} f(k) = \sum_{i=a}^{b} f(i)\).
Its name carries no information about the value of the sum.
\end{remark*}

\section*{Basic Properties of Summation}

\begin{remark*}[Basic properties of summation]
For functions \(f\) and \(g\) and constants \(c\),
\begin{align*}
\sum_{k=m}^{n} (f(k) + g(k))
&=
\sum_{k=m}^{n} f(k)
+
\sum_{k=m}^{n} g(k), \\[6pt]
\sum_{k=m}^{n} c\,f(k)
&=
c \sum_{k=m}^{n} f(k).
\end{align*}
These express the \emph{linearity} of summation.
\end{remark*}

\section*{Splitting Sums}

\begin{remark*}[Splitting sums]
A sum over an interval may be divided into smaller pieces:
\[
\sum_{k=a}^{c} f(k)
=
\sum_{k=a}^{b} f(k)
+
\sum_{k=b+1}^{c} f(k)
\qquad \text{whenever } a \le b < c.
\]
\end{remark*}

\section*{Changing the Index}

\begin{remark*}[Changing the index]
It is often useful to shift the summation index. For example,
\[
\sum_{k=0}^{n} f(k) = \sum_{j=1}^{n+1} f(j-1).
\]
Such substitutions frequently simplify expressions or align sums
with known formulas.
\end{remark*}

\section*{Empty Sums}

\begin{remark*}[Empty sums]
If the upper bound is smaller than the lower bound, the sum contains
no terms. By convention,
\[
\sum_{k=a}^{b} f(k) = 0 \qquad \text{when } a > b.
\]
This convention simplifies many algebraic identities.
\end{remark*}

\section*{Nested Summations}

\begin{remark*}[Nested summations]
Summations may be nested:
\[
\sum_{i=1}^{n} \sum_{j=1}^{m} a_{ij}
= \sum_{i=1}^{n}\bigl(a_{i1} + a_{i2} + \cdots + a_{im}\bigr).
\]
The inner sum is evaluated first. When the limits are independent,
the order of summation may be exchanged:
\[
\sum_{i=1}^{n}\sum_{j=1}^{m} a_{ij}
=
\sum_{j=1}^{m}\sum_{i=1}^{n} a_{ij}.
\]
\end{remark*}

\begin{remark*}[Analogy with integration]
Nested sums behave like iterated integrals: exchanging summation order
corresponds to Fubini's theorem in the discrete setting.
\end{remark*}

\section*{Common Examples}

\begin{remark*}[Common examples]
\[
\sum_{k=1}^{n} 1 = n, \qquad
\sum_{k=1}^{n} k = \frac{n(n+1)}{2}, \qquad
\sum_{k=1}^{n} k^2 = \frac{n(n+1)(2n+1)}{6}.
\]
These formulas appear throughout discrete mathematics and analysis.
\end{remark*}
